\documentclass[12pt, oneside]{article}

\usepackage{xparse}

\usepackage{geometry}
\geometry{letterpaper}
\usepackage[parfill]{parskip}
\usepackage{float}

% Math packages
\usepackage{amssymb}
\usepackage{amsmath}
\usepackage{amsthm}
\usepackage{mathalpha}

% Tables
\usepackage{booktabs}
\usepackage{afterpage}

% Hyperlinks
\usepackage{hyperref}

% References
\numberwithin{equation}{section}
\usepackage[numbers]{natbib}

% TikZ (optional for later)
\usepackage{tikz}
\usetikzlibrary{shapes,arrows,chains}

% Theorem environments
\theoremstyle{definition}
\newtheorem{definition}{Definition}[section]
\newtheorem{proposition}{Proposition}[section]
\newtheorem{axiom}{Axiom}[section]

\theoremstyle{plain}
\newtheorem{theorem}{Theorem}[section]
\newtheorem{lemma}[theorem]{Lemma}
\newtheorem{corollary}[theorem]{Corollary}
\newtheorem{example}[theorem]{Example}

\theoremstyle{remark}
\newtheorem*{remark}{Remark}

\title{Affine Constraint Dynamics and Density-One Descent in the Collatz Map}
\author{Wayne Brassem}

\begin{document}

% Custom commands
\NewDocumentCommand{\setN}{}{\mathbb{N}}
\NewDocumentCommand{\setZ}{}{\mathbb{Z}}
\NewDocumentCommand{\setQ}{}{\mathbb{Q}}

% Document commands which provide hyperlinks to OEIS sequences referenced herein
\NewDocumentCommand{\OEIS}{m}{\href{https://oeis.org/#1}{#1}}

\maketitle

\begin{abstract}

We study the Collatz map via the fully accelerated transformation
\[
f(x)=\frac{3x+1}{2^{v_2(3x+1)}}.
\]

Finite trajectories are encoded by division words, which record the
2-adic valuations encountered along an orbit. Each such word induces
an affine map whose admissibility imposes a congruence constraint on
initial values. This yields a symbolic–affine framework in which
admissible words correspond to residue classes modulo powers of two.

We show that these constraints refine exponentially: a division word
of length $m$ restricts initial values to a residue class modulo a
power of $2$ that is comparable to, but strictly larger than, $3^m$,
while the number of admissible words grows like $\rho^m$ with $\rho < 3$.

This creates a competition between combinatorial growth and arithmetic
refinement: while the number of admissible words grows like $\rho^m$
with $\rho < 3$, each word imposes a congruence constraint modulo a
power of $2$ that grows strictly faster than $3^m$. This imbalance
forces exponential sparsification of admissible initial conditions.

As a consequence, the set of integers that fail to admit a finite
descent prefix has natural density zero. In particular, almost all
integers eventually decrease under iteration of the Collatz map.

\end{abstract}

\newpage
\tableofcontents

% ============================================================
\newpage
\section{Introduction}

The $3x+1$ problem, also known as the Collatz conjecture~\cite{Collatz1937}, concerns
the behavior of the map
\[
T(x) =
\begin{cases}
\displaystyle
    \phantom{3x} \frac{x}{2} & \text {if } x \text{ is even}, \\[8pt]
\displaystyle
    3x+1,                    & \text{if $x$ is odd}.
\end{cases}
\]
on the positive integers. It asserts that every positive integer eventually reaches $1$
under iteration of $T$. Despite its elementary formulation, the conjecture has resisted
proof for decades~\cite{Lagarias2010,Terras1976}. Extensive computational verification
confirms convergence for very large ranges~\cite{OliveiraESilva2010}, but does not
reveal the global structure governing trajectory behavior.

A major advance by Tao~\cite{Tao2019} showed that almost all integers, in the sense of
logarithmic density, eventually decrease under iteration of the Collatz map. This result
is obtained via probabilistic averaging rather than an explicit structural decomposition
of the integers.

The present work develops a deterministic, arithmetic framework for analyzing typical
Collatz behavior based on a symbolic–affine encoding of trajectories. Using this approach,
we encode the dynamics symbolically rather than studying individual trajectories with
the fully accelerated map
\[
f(x)=\frac{3x+1}{2^{v_2(3x+1)}},
\]
which maps odd integers to odd integers by compressing each odd step into a single
transformation.

Each trajectory under $f$ is represented by a finite \emph{division word}, recording the
number of factors of $2$ removed (the 2-adic valuation) at each iterate.

These words admit an exact affine representation obtained by symbolic composition of the
accelerated map: each division word induces a map of the form
\[
F_\omega(x) = \frac{3^m x + c(\omega)}{2^{D(\omega)}},
\]
and determines a corresponding congruence class modulo $2^{D(\omega)}$. Thus division words
determine canonical arithmetic \emph{constraint classes} (residue classes) in the integers,
once realizability is imposed.

This correspondence allows the dynamics to be studied at the level of residue classes
rather than individual orbits. Convergent division words—those for which the associated
affine map is contractive—identify families of integers that decrease after a fixed
number of steps. The global problem is thereby reduced to understanding how these
families are distributed and how their densities accumulate.

Two competing effects govern this distribution. On one hand, the number of admissible
convergent division words grows exponentially with length. On the other hand, each word
imposes a congruence constraint whose modulus grows exponentially with the total 2-adic
valuation. Using structural constraints on admissible division words, encoded by an
increment grammar, we show that the combinatorial growth rate is strictly subcritical:
the number of admissible words grows like $\rho^m$ for some $\rho < 3$, while the moduli
scale like powers of $2$ satisfying $3^m < 2^{D(\omega)} < 2\cdot 3^m$.

This transforms the Collatz problem into a competition between symbolic growth and
arithmetic refinement, with convergence governed by their relative exponential rates.

This imbalance leads to \emph{exponential sparsification}: the proportion of integers
admitting a convergent division word of length $m$ decays exponentially in $m$.
Summing over all lengths then shows that the complement—integers that avoid all such
descent patterns—has natural density zero.

Consequently, the set of integers whose trajectory under iteration of $f$ eventually
decreases has natural density one.

While this does not exclude the existence of exceptional trajectories, any such
exceptions must lie within a set of zero density and cannot arise from any scalable
arithmetic structure generated by the symbolic dynamics.

\paragraph{Structure of the Paper.}
The argument proceeds as follows:
\begin{itemize}
\item Section~\ref{sec:division-counts} introduces division words and their total division counts.
\item Section~\ref{sec:affine-representation} develops the affine representation of division words and their congruence structure.
\item Section~\ref{sec:affine-grammar} establishes admissibility and minimal valuation constraints.
\item Section~\ref{sec:symbolic-growth} controls the growth of admissible words via the increment grammar.
\item Section~\ref{sec:affine-constraint-fibers} identifies the residue class structure induced by division words.
\item Section~\ref{sec:exponential-sparsification} proves exponential sparsification of admissible sets.
\item Section~\ref{sec:density-one-descent} establishes density-one descent.
\end{itemize}

% ============================================================
\section{Symbolic Encoding of Collatz Trajectories}
\label{sec:division-counts}

% ------------------------------------------------------------
\subsection{Division Words}

\begin{definition}[Division Word]
A division word of length $m \ge 0$ is a sequence
\[
\omega = (d_0, \dots, d_{m-1}),
\]
where each $d_i \in \mathbb{N}$ is intended to represent the number of factors of $2$
removed at the $i$-th iterate of the accelerated map
\[
f(x) = \frac{3x+1}{2^{v_2(3x+1)}}.
\]
\end{definition}

\begin{remark}
Since $f$ maps odd integers to odd integers, each $3f^i(x)+1$ is even,
so $d_i \ge 1$ for all $i$.
\end{remark}

\begin{definition}[Realization]
An odd integer $x \in \mathbb{N}$ realizes $\omega$ if
\[
d_i = v_2(3f^i(x)+1), \quad 0 \le i < m.
\]
\end{definition}

% ------------------------------------------------------------
\subsection{Total Division Count}

\begin{definition}
The total division count of a word $\omega$ is
\[
D(\omega) := \sum_{i=0}^{m-1} d_i.
\]
\end{definition}

\begin{remark}
The quantity $D(\omega)$ measures the total dyadic contraction
associated with the symbolic path $\omega$.
\end{remark}

% ============================================================
\section{Affine Representation of Division Words}
\label{sec:affine-representation}

\begin{definition}
Each division word $\omega$ of length $m$ induces an associated affine map
\[
F_\omega(x) = \frac{3^m x + c(\omega)}{2^{D(\omega)}},
\]
obtained by symbolic composition of the accelerated map.
\end{definition}

\begin{proposition}
The map $\omega \mapsto F_\omega$ defines a semigroup homomorphism from
concatenation of division words (read left to right) to composition of affine maps.
\end{proposition}

% ------------------------------------------------------------
\subsection{Example}

Consider convergent segments with $m=4$ iterates of $f$.  
The minimal total division count for $m=4$ is
\[
A020914(4) = 7.
\]

One admissible division word attaining this minimum is
\[
\omega = (1,1,1,4),
\]
for which
\[
D(\omega) = 1+1+1+4 = 7.
\]

Composing the accelerated map according to $\omega$ yields
\[
f_\omega(x)=\frac{3^4 x + c(\omega)}{2^7}
           = \frac{81x + 65}{128},
\]
so the canonical affine constant for this word is
\[
c(\omega)=65.
\]

The integrality condition
\[
81x + 65 \equiv 0 \pmod{128}
\]
determines a unique residue class
\[
x \equiv 15 \pmod{128}.
\]
This congruence defines a constraint fiber with modulus $2^7$.

\medskip

\noindent
\textbf{Interpretation.}
The word $\omega = (1,1,1,4)$ means:

\begin{itemize}
    \item At each of the first three iterates of $f$, $3x+1$ contains exactly one factor of $2$.
    \item At the fourth iterate of $f$, $3x+1$ contains four factors of $2$.
\end{itemize}

Under the accelerated map
\[
f(x)=\frac{3x+1}{2^{v_2(3x+1)}},
\]
all of these divisions occur within a single iterate. Thus $d_3=4$ means that
the fourth iterate divides by $2^4$ in that step. The constant $c(\omega)$ is
fixed entirely by this symbolic structure and ensures that precisely the
integers in the residue class $x \equiv 15 \pmod{128}$ produce integer iterates
under the affine form.

This example illustrates that the behavior of trajectories is governed
entirely by the symbolic data encoded in $\omega$ and its associated
affine map. The initial values that realize a given word are precisely
those satisfying the induced congruence constraint, independent of any
particular choice of representative.

\medskip

\noindent
In general, not every division word arises from an actual trajectory;
the condition for realizability will be formulated later as an explicit
congruence constraint, defining the class of \emph{admissible} words.

% ============================================================
\section{Affine Symbolic Admissibility Grammar}
\label{sec:affine-grammar}

\begin{remark}[Division Grammar as a Dynamical Encoding]
Recall from Section~\ref{sec:division-counts} that a \emph{division word}
$\omega=(d_0,\dots,d_{m-1})$ records the exact 2-adic valuation
$d_i = v_2(3 x_i + 1)$, occurring at each iterate of the fully accelerated
Collatz map.

Unlike a purely combinatorial encoding, division words retain essential
dynamical information. Each division word $\omega=(d_0,\dots,d_{m-1})$
determines the sequence of multiplicative factors of the linear coefficient
\[
r_i = \frac{3}{2^{d_i}},
\]
so that the contractive or expansive contribution of each iterate is
preserved in the linear coefficient (with the affine term handled separately
via $c(\omega)$). This structure enables quantitative control of admissible
trajectories via affine growth and to formulate sufficient conditions for
descent, even though individual trajectories are collapsed under the encoding.
\end{remark}

% ------------------------------------------------------------
\subsection{Explicit Form of the Affine Constant}
\begin{lemma}[Explicit Form of the Affine Constant]
Let $\omega = (d_0,\dots,d_{m-1})$ be a division word. Then the affine map
\[
F_\omega(x) = \frac{3^m x + c(\omega)}{2^{D(\omega)}}
\]
has the explicit form
\[
c(\omega)
=
\sum_{j=0}^{m-1}
3^{m-1-j}\cdot 2^{\sum_{i=0}^{j-1} d_i},
\]
where empty sums are defined as $0$.
\end{lemma}

% ------------------------------------------------------------
\subsection{Injectivity of the Affine Encoding}
\begin{lemma}[Injectivity of the Affine Encoding at Fixed Length]
\label{lem:injectivity}
Let $\omega = (d_0,\dots,d_{m-1})$ and $\omega' = (d'_0,\dots,d'_{m-1})$
be division words of the same length $m$. Let
\[
c(\omega)
=
\sum_{j=0}^{m-1}
3^{m-1-j}\cdot 2^{D_j},
\quad
D_j = \sum_{i=0}^{j-1} d_i,
\]
and define $c(\omega')$ analogously.

If $\omega \ne \omega'$, then
\[
c(\omega) \ne c(\omega').
\]

Equivalently, the map $\omega \mapsto c(\omega)$ is injective on admissible
division words of fixed length $m$.
\end{lemma}

\begin{proof}
Suppose $\omega \ne \omega'$, and let
\[
k = \min\{\, j : d_j \ne d'_j \,\}
\]
be the first index at which the two words differ. Then for all $j < k$,
we have $D_j = D'_j$, while $D_k \ne D'_k$.

Consider the difference
\[
\Delta := c(\omega) - c(\omega')
=
\sum_{j=0}^{m-1}
3^{m-1-j}\bigl(2^{D_j} - 2^{D'_j}\bigr).
\]
The terms with $j<k$ cancel, so
\[
\Delta
=
3^{m-1-k}\bigl(2^{D_k} - 2^{D'_k}\bigr)
+
\sum_{j=k+1}^{m-1}
3^{m-1-j}\bigl(2^{D_j} - 2^{D'_j}\bigr).
\]

Let $\delta_k = \min(D_k, D'_k)$. Without loss of generality,
assume $D_k > D'_k$. Then
\[
2^{D_k} - 2^{D'_k}
=
2^{D'_k}(2^{D_k - D'_k} - 1),
\]
where $2^{D_k - D'_k} - 1$ is odd. Hence
\[
v_2(2^{D_k} - 2^{D'_k}) = D'_k = \delta_k.
\]

For $j>k$, both $D_j$ and $D'_j$ strictly exceed $\delta_k$, since the
prefix sums increase by positive integers. Therefore
\[
v_2\!\left(3^{m-1-j}\bigl(2^{D_j} - 2^{D'_j}\bigr)\right)
>
\delta_k
\quad \text{for all } j>k.
\]

It follows that $\Delta$ is a sum of one term with 2-adic valuation
exactly $\delta_k$ and remaining terms of strictly higher 2-adic valuation,
so no cancellation at the lowest 2-adic order is possible. Hence
\[
v_2(\Delta) = \delta_k,
\]
and in particular $\Delta \ne 0$.

Therefore $c(\omega) \ne c(\omega')$, completing the proof.
\end{proof}

This implies distinct division words induce distinct affine constraints.

% ------------------------------------------------------------
\subsection{Stepwise Realizability}

\begin{theorem}[Stepwise Realizability of Division Words]
\label{thm:stepwise-realizability}
Let $\omega = (d_0,\dots,d_{m-1})$ be a division word, and define $c(\omega)$ by
\[
c(\omega)
=
\sum_{j=0}^{m-1}
3^{m-1-j}\cdot 2^{\sum_{i=0}^{j-1} d_i}.
\]

Then the identity
\[
2^{D(\omega)} x_m = 3^m x_0 + c(\omega)
\]
holds if and only if the sequence $(x_0,\dots,x_m)$ satisfies the accelerated
Collatz relations
\[
x_{i+1} = \frac{3x_i+1}{2^{d_i}}, \quad 0 \le i < m,
\]
with each intermediate step integral.

In particular, the congruence
\[
3^m x_0 + c(\omega) \equiv 0 \pmod{2^{D(\omega)}}
\]
is equivalent to the existence of a trajectory realizing $\omega$.
\end{theorem}

In particular, the congruence ensures that each intermediate iterate
$x_i$ is an integer, so realizability is enforced at every step.

\begin{remark}[Structural Interpretation]
Each term in $c(\omega)$ corresponds to the additive contribution of a single
``+1'' in the Collatz iteration, propagated forward by subsequent multiplications
by $3$ and backward-scaled by preceding divisions by $2$.
\end{remark}

% ------------------------------------------------------------
\subsection{Admissibility of Division Words}

\begin{definition}[Admissibility]
\label{def:admissible-division-word}
A division word $\omega$ is admissible if the congruence
\[
3^m x + c(\omega) \equiv 0 \pmod{2^{D(\omega)}}
\]
has a solution.
\end{definition}
\begin{remark}
Admissibility is therefore an arithmetic condition on the word $\omega$
itself, expressed as a divisibility constraint on the associated affine
map. When this condition holds, the set of all integers $x$ for which
$F_\omega(x)$ is integral defines a unique congruence constraint modulo
$2^{D(\omega)}$, corresponding to the residue class induced by $\omega$.
\end{remark}

% ------------------------------------------------------------
\subsection{Minimal Division Count}

\begin{lemma}[Minimal Division Count]
For any admissible division word $\omega=(d_0,\dots,d_{m-1})$ of length $m$,
\[
d_0 + \cdots + d_{m-1} \;\ge\; A020914(m),
\]
where
\[
A020914(m) = \lceil m \log_2(3) \rceil,
\]
the minimal integer $k$ such that $2^k > 3^m$.

Equality holds if and only if $\omega$ achieves the minimal possible total
division count among admissible words of length $m$.
\end{lemma}

% ------------------------------------------------------------
\subsection{Convergent Division Word}

\begin{definition}[Convergent Division Word]
An admissible division word $\omega = (d_0,\dots,d_{m-1})$ is said to be
\emph{convergent} (in the linear sense) if
\[
\frac{3^m}{2^{d_0+\cdots+d_{m-1}}} < 1.
\]
\end{definition}

% ============================================================
\section{Symbolic Growth and Spectral Bounds}
\label{sec:symbolic-growth}

The combinatorial complexity of admissible division words governs the number of
affine constraint fibers at each depth. This growth is encoded by the sequence
\OEIS{A186009}, which counts convergent division words and hence, by injectivity
of the affine encoding, the associated admissible residue classes by stopping time.

% ------------------------------------------------------------
\subsection{Combinatorial growth model}

The sequence \OEIS{A186009} arises from a horizontal-sum recurrence of Pascal type,
augmented by boundary duplication events governed by the increment grammar
\OEIS{A022921}. This grammar induces a finite symbolic system on two local operations:
single increments ($s$) and double increments ($d$), whose admissible
concatenations generate all configurations counted by \OEIS{A186009}.

The resulting dynamics decompose into a finite set of admissible local patterns,
whose concatenations generate all admissible configurations. The key structural
feature of this grammar is that locally expansive configurations cannot occur
arbitrarily often: they are separated by mandatory contractive patterns.

As a consequence, the total number of admissible configurations exhibits
\emph{submultiplicative growth}, and admits a well-defined exponential growth
rate in the sense of a limsup.

% ------------------------------------------------------------
\subsection{Subcritical growth bound}

\begin{theorem}[Subcritical symbolic growth]
\label{thm:subcritical-growth}
Let $\widetilde A(i) = A186009(i+1)$ denote the number of admissible configurations
of stopping time $i$ (i.e., convergent admissible division words of length $i$).
Then there exists a constant $\rho < 3$ such that
\[
\widetilde A(i) \le C \rho^i.
\]
In particular,
\[
\limsup_{i\to\infty} \widetilde A(i)^{1/i} \le \rho < 3.
\]
\end{theorem}

\begin{remark}
The constant $\rho$ arises from the maximal spectral growth of admissible
cycles in the increment grammar. A complete derivation is given in
Appendix~\ref{app:spine-growth}.
\end{remark}

% ------------------------------------------------------------
\subsection{Dyadic comparison}

Let $B(i) = A020914(i) = \lceil i \log_2 3 \rceil$. Thus $2^{B(i)}$ is the minimal
power of $2$ exceeding $3^i$, and
\[
3^i < 2^{B(i)} < 2 \cdot 3^i,
\]
so the dyadic scaling associated with admissible constraint fibers grows
asymptotically like $3^i$.

Indeed, combining $\widetilde A(i) \le C\rho^i$ with $2^{B(i)} > 3^i$ gives
\[
\frac{\widetilde A(i)}{2^{B(i)}} \le C\left(\frac{\rho}{3}\right)^i \to 0
 \quad \text{as } i \to \infty,
\]
since $\rho < 3$.

Thus the exponential growth of admissible configurations is strictly dominated
by the exponential growth of their associated congruence moduli.

% ------------------------------------------------------------
\subsection{Spectral bracketing}

The symbolic system admits explicit lower and upper envelope models.

\paragraph{Lower envelope.}
The restricted grammar underlying \OEIS{A047749}, which forbids consecutive
doubling events, yields an exactly solvable comparison sequence with
exponential growth rate
\[
\lambda_{\mathrm{low}} = \frac{3\sqrt{3}}{2}.
\]
This construction is detailed in Appendix~\ref{app:A047749-lower-bound}.

\paragraph{Upper envelope.}
The full admissible grammar of \OEIS{A186009} admits a maximal growth rate
\[
\lambda_{\mathrm{high}} = \rho < 3,
\]
arising from the dominant cycle structure of the increment grammar
(Appendix~\ref{app:spine-growth}).

\medskip

\noindent
These bounds yield a spectral bracketing:
\[
\lambda_{\mathrm{low}}
\;\le\;
\limsup_{i\to\infty} \widetilde A(i)^{1/i}
\;\le\;
\lambda_{\mathrm{high}}.
\]

\begin{remark}[Reduction to a Single Growth Parameter]
All combinatorial complexity of admissible division words is absorbed into
the single parameter $\rho < 3$, representing the maximal exponential growth
rate of the symbolic system. Subsequent arguments depend only on this
subcritical growth bound, and not on the detailed structure of the enumeration.
\end{remark}

\noindent
From this point forward, the detailed combinatorics of admissible division
words enter only through the growth bound $\rho < 3$.

\subsection*{Interpretation}

The admissible symbolic dynamics are therefore confined to a strictly
subcritical regime: although the number of configurations grows exponentially,
its rate is bounded away from the dyadic scaling rate $3$.

This establishes a fundamental imbalance: combinatorial expansion of admissible
division words grows strictly slower than the dyadic refinement imposed by their
congruence constraints.

This imbalance is the mechanism driving exponential sparsification in the next section.

% ============================================================
\section{Affine Constraint Fibers}
\label{sec:affine-constraint-fibers}

% ------------------------------------------------------------
\subsection{Congruence Stability of Division Words}

\begin{theorem}[Congruence Stability of Division Words]
Let $\omega$ be admissible and let $r_\omega$ satisfy
\[
r_\omega \equiv -c(\omega)\cdot (3^m)^{-1} \pmod{2^{D(\omega)}}.
\]
Then for all $x \equiv r_\omega \pmod{2^{D(\omega)}}$, the first $m$ iterations
of the accelerated Collatz map realize exactly the division word $\omega$,
by Theorem~\ref{thm:stepwise-realizability}.
\end{theorem}

% ------------------------------------------------------------
\subsection{Residue Class Representation of Division Words}

\begin{theorem}[Affine Constraint Representation]
Let $\omega$ be an admissible division word of length $m$.

Then there exists a unique residue $r_\omega \pmod{2^{D(\omega)}}$ such that
\[
x \text{ realizes } \omega \iff x \equiv r_\omega \pmod{2^{D(\omega)}}.
\]

Thus admissible division words correspond to affine congruence constraints
of modulus $2^{D(\omega)}$.
\end{theorem}

\begin{proof}
From the affine representation,
\[
F_\omega(x) = \frac{3^m x + c(\omega)}{2^{D(\omega)}},
\]
integrality is equivalent to
\[
3^m x + c(\omega) \equiv 0 \pmod{2^{D(\omega)}}.
\]

Since $3^m$ is odd, it is invertible modulo $2^{D(\omega)}$, yielding
\[
x \equiv -c(\omega) \cdot (3^m)^{-1} \pmod{2^{D(\omega)}}.
\]

This determines a unique residue modulo $2^{D(\omega)}$.

Conversely, any $x$ in this class satisfies the integrality condition.
By Theorem~\ref{thm:stepwise-realizability}, this is equivalent to the
existence of a trajectory whose successive iterates satisfy
\[
x_{i+1} = \frac{3x_i+1}{2^{d_i}}, \quad 0 \le i < m,
\]
and hence realize the division word $\omega$.

Thus admissible division words are in one-to-one correspondence with affine
congruence constraints modulo powers of two.
\end{proof}

% ============================================================
\section{Exponential Sparsification}
\label{sec:exponential-sparsification}

\begin{lemma}[Constraint Refinement]
Each admissible division word $\omega$ imposes a congruence constraint
modulo $2^{D(\omega)}$, where $D(\omega) \ge A020914(m)$.
\end{lemma}

\begin{lemma}[Subcritical Symbolic Growth]
The number of admissible division words of length $m$ satisfies
\[
\#\{\omega : |\omega|=m\} = O(\rho^m)
\quad \text{for some } \rho < 3.
\]
\end{lemma}

\begin{proof}
This follows from Theorem~\ref{thm:subcritical-growth}, which bounds the growth
of admissible configurations generated by the increment grammar.
\end{proof}

\begin{theorem}[Exponential Sparsification]
The set of integers admitting a division word of length $m$ occupies at most
a proportion on the order of
\[
\frac{\#\{\text{admissible words of length } m\}}{2^{A020914(m)}}.
\]
In particular, this proportion decays exponentially in $m$.
\end{theorem}

\begin{lemma}[Exponential Tail Decay]
Let $E_m$ denote the set of integers that do not admit a descent prefix
of length at most $m$.

Then $\mu(E_m)$ decays exponentially in $m$.
\end{lemma}

\begin{proof}
An integer fails to admit a descent prefix of length at most $m$ only if it avoids
all admissible convergent division words of length $k \le m$.

Each such word $\omega$ of length $k$ imposes a congruence constraint modulo
$2^{D(\omega)}$, with $D(\omega) \ge A020914(k)$.

For each fixed $k$, the proportion of integers admitting a convergent division
word of length $k$ is bounded by
\[
\sum_{|\omega|=k} \frac{1}{2^{D(\omega)}}
\;\le\;
C\left(\frac{\rho}{3}\right)^k,
\]
using $D(\omega) \ge A020914(k)$ and $\rho < 3$.

Summing over $k \le m$ yields
\[
\sum_{k=1}^{m} \sum_{|\omega|=k} \frac{1}{2^{D(\omega)}}
\;\le\;
C \sum_{k=1}^{m} \left(\frac{\rho}{3}\right)^k,
\]
which converges exponentially as $m \to \infty$.

Thus the total proportion of integers admitting a descent prefix of length at most $m$
approaches $1$ exponentially fast, and the complement $E_m$ decays exponentially.
\end{proof}

% ============================================================
\section{Density-One Descent}
\label{sec:density-one-descent}

\begin{lemma}[Exhaustion by Descent Prefixes]
Let $S_m$ denote the set of integers that admit no convergent division word
of length $\le m$. Then
\[
S_{m+1} \subseteq S_m,
\quad\text{and}\quad
\bigcap_{m=1}^\infty S_m
\]
is precisely the set of integers that admit no finite descent prefix.
\end{lemma}

\begin{theorem}[Density-One Descent]
The set of integers that do not admit any finite convergent division word
has natural density zero.
\end{theorem}

\begin{proof}[Proof of Theorem]
For each $m$, the complement of $S_m$ is the union of all residue classes
corresponding to convergent division words of length $\le m$.

Each convergent division word $\omega$ of length $k \le m$ defines a residue
class of density $2^{-D(\omega)} \le 2^{-B(k)}$, and by subadditivity of
natural density (i.e., a union bound), the total measure is bounded above by
the sum of their individual densities.

Since $2^{B(k)} > 3^k$ and $\widetilde A(k) \le C \rho^k$ with $\rho < 3$,
thus the total contribution over all lengths $k \le m$ satisfies
\[
\sum_{k=1}^{m} \sum_{|\omega|=k} \frac{1}{2^{D(\omega)}}
\;\le\;
C \sum_{k=1}^{m} \left(\frac{\rho}{3}\right)^k,
\]
which decays exponentially.

Summing these exponentially decaying contributions over all lengths $\le m$
shows that the total density of integers admitting a convergent division word
of length at most $m$ tends to $1$ as $m \to \infty$.

Therefore
\[
\lim_{m\to\infty} \mu(S_m^c) = 1
\quad\text{and hence}\quad
\mu\!\left(\bigcap_{m=1}^\infty S_m\right) = 0,
\]
which proves the theorem.
\end{proof}

\begin{corollary}
The set of integers whose trajectory under $f$ eventually decreases
has natural density one.
\end{corollary}

\begin{remark}
Density emerges as a consequence of affine growth imbalance.
\end{remark}

% ------------------------------------------------------------
\subsection{Hypothetical Zero-Density Trajectories}

Any trajectory avoiding all convergent division words must lie outside
the union of convergent constraint classes whose cumulative density
approaches one. Such trajectories, if they exist, are therefore confined
to a set of natural density zero.

In particular, no exceptional behavior can arise from any scalable
arithmetic structure generated by the symbolic dynamics. Any such
exception would be dynamically isolated within a sparse residual set,
rather than supported by a hierarchy of positive-density constraint
classes.

% ============================================================
\section{Discussion and Outlook}

The preceding results establish a density-theoretic dominance of
dyadic constraint refinement over the combinatorial growth of admissible
division words.

Several aspects lie outside the present framework, including the
fine structure of atypical trajectories within the zero-density residual
set and possible refinements of the increment grammar toward sharper
spectral bounds.

The affine–symbolic framework developed here is compatible with
more refined local analyses, particularly in regimes where the
minimal valuation structure becomes rigid.

\newpage
\appendix
\section{Division Words, Grammar Structure, and Spectral Growth}
\label{app:spine}

This appendix collects the structural framework underlying the division-word encoding,
its affine cylinder representation, the admissibility grammar governed by \OEIS{A020914},
and the resulting growth theory for \OEIS{A186009}. The presentation is organized in three layers:
encoding, grammar, and growth.

% ============================================================
\subsection{Affine Encoding of Division Words}
\label{app:spine-affine}

A division word of length $m \ge 1$ is a tuple
\[
\omega = (d_0, d_1, \dots, d_{m-1}), \qquad d_i \in \mathbb{N},
\]
with total exponent
\[
D(\omega) = \sum_{i=0}^{m-1} d_i.
\]

Each word induces an affine transformation
\[
F_\omega(x) = \frac{3^m x + c(\omega)}{2^{D(\omega)}}.
\]

The constant $c(\omega)$ is defined modulo $2^{D(\omega)}$ by
\[
3^m x + c(\omega) \equiv 0 \pmod{2^{D(\omega)}},
\qquad 0 \le c(\omega) < 2^{D(\omega)}.
\]

Let
\[
S_j = \sum_{i=0}^{j-1} d_i, \quad S_0 = 0.
\]

Then
\[
c(\omega)
=
\sum_{j=0}^{m-1}
3^{m-1-j} 2^{S_j}.
\]

Since $3^m$ is invertible modulo $2^{D(\omega)}$, admissible inputs satisfy
\[
x \equiv -c(\omega)(3^m)^{-1} \pmod{2^{D(\omega)}}.
\]

Thus each word defines a cylinder
\[
\mathcal{C}_\omega = \{ r_\omega + k2^{D(\omega)} : k \ge 0 \}.
\]

% ============================================================
\subsection{Grammar Structure and Enumeration}
\label{app:spine-grammar}

Let
\[
H(i) = \sum_{j=0}^i d_j,
\qquad A(i) = \text{\OEIS{A020914}}(i).
\]

A division word is admissible iff
\[
H(i) < A(i+1) \quad (i < m-1),
\qquad
H(m-1) = A(m).
\]

Admissible words are generated by bounded backtracking over
\[
d_i \in \mathbb{N}_{\ge 1}
\]
subject to the prefix constraints.

For $m=5$, with $A(1{:}5) = (2,4,5,7,8)$, there are 7 admissible words,
matching $\text{\OEIS{A186009}}(6)=7$.

Each word corresponds bijectively to a cylinder:
\[
\text{division word} \;\leftrightarrow\; \text{cylinder} \;\leftrightarrow\; \text{affine progression}.
\]

% ============================================================
\subsection{Horizontal-Sum Dynamics and Growth}
\label{app:spine-growth}

The sequence \OEIS{A186009} arises from
\[
a_i(n) = a_i(n-1) + a_{i-1}(n),
\]
with boundary modifications governed by A022921.

This induces local patterns:
\[
(s,d), \quad (d,d), \quad (s,d,s), \quad (d,d,s).
\]

The multiplicative factors associated with each admissible local pattern are:
\[
\begin{array}{c|c}
\text{Pattern} & \text{Factor} \\
\hline
(s,d) & 3 \\
(d,d) & 123/33 \\
(s,d,s) & 341/131 \\
(d,d,s) & 329/123
\end{array}
\]

By monotonicity of the horizontal-sum recurrence, no admissible realization
of a given pattern can exceed its extremal value, so these factors provide
uniform upper bounds across all admissible configurations.

These values arise from extremal realizations of the horizontal-sum dynamics;
a detailed derivation is provided in Appendix~\ref{app:local-growth-compressed}, 
see in particular Section~\ref{subsec:extremal-cycle}.

Admissible sequences decompose into cycles. The extremal 5-cycle gives
\[
\rho_5 =
\left(
3^2 \cdot \frac{123}{33} \cdot \frac{341}{131} \cdot \frac{329}{123}
\right)^{1/5}
< 3.
\]

Thus
\[
\limsup_{n\to\infty} A(n)^{1/n} \le \rho_5.
\]

Since $2^{A020914(n)} \sim 3^n$,
\[
\frac{A(n)}{2^{A020914(n)}} \to 0.
\]

% ------------------------------------------------------------
\subsubsection*{Lower spectral envelope (A047749)}
\label{app:A047749-lower-bound}

\OEIS{A047749} corresponds to the restricted grammar forbidding consecutive doubling.

Pointwise:
\[
A047749(n) \le A186009(n+1).
\]

Closed form:
\[
a(2m) = \frac{1}{2m+1}\binom{3m}{m}, \quad
a(2m+1) = \frac{1}{2m+1}\binom{3m+1}{m+1}.
\]

Asymptotic growth:
\[
a(n) \sim C\left(\frac{3\sqrt{3}}{2}\right)^n.
\]

Hence lower rate:
\[
\lambda_{\min} = \frac{3\sqrt{3}}{2}.
\]

% ------------------------------------------------------------
\subsubsection*{Spectral sandwich}

\[
\left(\frac{3\sqrt{3}}{2}\right)^n
\;\lesssim\;
A(n)
\;\ll\;
2^{A020914(n)}.
\]

\[
\text{division words}
\Rightarrow
\text{cylinders}
\Rightarrow
\text{grammar}
\Rightarrow
\text{cycles}
\Rightarrow
\text{spectral bounds}.
\]

% ============================================================
\section{Local Growth Modes and Spectral Reduction}
\label{app:local-growth-compressed}

\subsection{Finite Grammar of Local Transitions}

The increment structure induced by \OEIS{A022921} admits a finite set of
admissible local configurations:
\[
\mathcal{P} = \{ (s,d), (d,d), (s,d,s), (d,d,s) \}.
\]

Each configuration represents a maximal locally consistent valuation pattern
for the accelerated map dynamics.

\subsection{Induced Affine Growth Weights}

Each pattern $\pi \in \mathcal{P}$ induces an affine update of the form
\[
x \mapsto \frac{3^{\ell(\pi)} x + \beta_\pi}{2^{\gamma(\pi)}},
\]
where $\ell(\pi)$ and $\gamma(\pi)$ depend only on the local valuation structure.

We define the associated multiplicative growth factor
\[
r_\pi := \frac{3^{\ell(\pi)}}{2^{\gamma(\pi)}}.
\]

\subsection{Extremal Realizations}

The admissible grammar admits extremal realizations of each pattern
corresponding to saturating configurations of the horizontal-sum dynamics.
These realizations yield the maximal admissible growth factors:

\[
\begin{array}{c|c}
\text{Pattern} & r_\pi \\
\hline
(s,d) & 3 \\
(d,d) & 123 / 33 \\
(s,d,s) & 341 / 131 \\
(d,d,s) & 329 / 123
\end{array}
\]

These values arise from explicit finite saturations of the horizontal-sum
recurrence and serve as uniform upper bounds for all admissible instances
of each pattern.

\subsection{Spectral Construction}

Admissible infinite sequences correspond to walks in a finite weighted directed
graph whose edges are labeled by the local patterns $\pi \in \mathcal{P}$
and weighted by $r_\pi$.

The global growth rate is therefore given by the spectral radius of the
associated transfer operator.

In particular, admissible cycles correspond to closed walks in this graph,
and their asymptotic contribution is given by geometric means of the
associated edge weights.

\subsection{Extremal Cycle and Spectral Bound}
\label{subsec:extremal-cycle}

The spectral bound is obtained from a finite-cycle relaxation of the increment grammar
underlying \OEIS{A022921}. In this setting, admissible sequences correspond to
walks on a finite state graph whose edges carry multiplicative weights given by the
local growth factors identified in Section~\ref{app:spine-growth}.

Among all directed cycles in this grammar graph, the $5$-cycle
\[
\{ s, d, s, d, d \}
\]
is the $5$-cycle provides an extremal candidate for maximizing the geometric mean
of local growth factors. Its significance is not that it dominates all admissible
sequences exactly, but that it provides an \emph{upper-envelope relaxation}:
it is obtained by removing the combinatorial constraints that would otherwise
prevent indefinite repetition of locally expansive patterns.

In particular, indefinite repetition of this cycle is not admissible within the
full grammar, since it violates the prefix constraints encoded by
\OEIS{A020914}. However, replacing the constrained system by this unconstrained
cycle can only increase growth, and therefore yields a valid upper bound on the
spectral radius.

This gives the relaxed spectral estimate
\[
\rho :=
\left(
3^2 \cdot \frac{123}{33} \cdot \frac{341}{131} \cdot \frac{329}{123}
\right)^{1/5} \approx 2.976 < 3,
\]
which bounds from above the true growth rate of admissible sequences.

\subsection{Conclusion}

All admissible growth in the symbolic system is governed by a finite
weighted grammar whose spectral radius satisfies
\[
\rho < 3.
\]

This establishes a strict subcritical regime for symbolic expansion.

% ============================================================
\clearpage
\bibliographystyle{unsrtnat}
\bibliography{src/latex/references}

\end{document}